\section{Related Work}
\label{sec:related}

\textbf{Spark configuration tuning.} Spark tuning is a mature field. Methods include learned
performance models \cite{nguyen2018tuning}, tuning from a single profiling run
\cite{buchaca2020yoro}, low-overhead tuning of Spark SQL \cite{xin2022locat}, online tuning
under several goals \cite{li2023online}, stage-level tuning together with adaptive query
execution \cite{lyu2024optimizer}, and trading tuning overhead against flexibility
\cite{li2025loftune}. A survey \cite{herodotou2020survey} lists cost-aware, workload-adaptive
tuning as an open problem. In this work the goal is speed, and tuning cost is something to
reduce, not something to weigh against the number of future runs.

\textbf{Workload-aware optimization.} Workload features guide transfer learning across jobs
\cite{shen2023rover}, adaptive execution in lakehouses \cite{xue2024aqe}, studies of cloud
workloads \cite{parayil2025workload} and parameter recommendation \cite{bao2018learning}.
These systems use features to find a faster setting sooner; they do not ask whether
optimizing is worth it at all.

\textbf{Compaction and data layout.} AutoComp \cite{gruenheid2025autocomp} automates
compaction in production data lakes and balances its benefit against its cost; it is the
closest work on compaction. Practitioner guides describe compaction for Iceberg and Delta Lake
\cite{goswami2023iceberg,delta2025compaction}, LST-Bench \cite{camacho2024lstbench} measures
table maintenance over time, and other work applies reinforcement learning to compaction
\cite{venkatesan2025autocompactionai} or improves small-file merging in HDFS
\cite{chettaoui2025merging}. PTO \cite{meduri2026pto} predicts with learned models whether a
layout change will help before applying it, AutoLiquid \cite{nakandala2026autoliquid}
automates layout in the Databricks lakehouse, and LayoutPilot \cite{liu2026layoutpilot} is a
physical-design advisor for lakehouse tables. None of them frames the choice as a break-even
over the future runs of a recurring pipeline.

\textbf{Cost and break-even.} Fekry \emph{et al.} \cite{fekry2020tune,fekry2020tuneful} show
that whether tuning is worthwhile depends on how often a job runs, and spread tuning cost over
future runs; we build on this idea. OptEx \cite{sidhanta2016optex} models cost and time to pick
a Spark cluster size under a deadline for a single job. A recent preprint
\cite{chang2026crosslayer} uses a measured break-even point for hardware-offload decisions in
databases. CherryPick \cite{alipourfard2017cherrypick}, Micky \cite{hsu2018micky}, PARIS
\cite{yadwadkar2017paris} and Smartpick \cite{mohapatra2023smartpick} choose cloud resources
for analytics jobs; they decide \emph{what to rent}, while we decide whether to optimize a
pipeline on resources already rented.

\textbf{Our position.} Prior systems automate one kind of optimization, often with learned
models and at production scale. We test one explicit break-even rule across an
execution-level and a storage-level action in a controlled setting.
